\clearpage
\VerdierPhysicalPageAnchor{0085}%
\VerdierSemanticAnchor{ch1-3-1-4-homotopy-definition-continuation}
是 $\operatorname{tr}_1(C)$ 的两个态射。连接 $(g,h,a)$ 与
$(g',h',a')$ 的一个\emph{同伦}是三元组 $(b,b',\alpha)$，其中
$b\colon g\rightsquigarrow g'$ 与 $b'\colon h\rightsquigarrow h'$ 是同伦，而

\VerdierSemanticAnchor{ch1-3-1-4-second-order-condition}
\begin{equation*}
  \alpha\colon f'\star b+a'\rightsquigarrow b'\star f+a
\end{equation*}
是二阶同伦 (2.5.12)。

\VerdierSemanticAnchor{ch1-3-1-4-homotopic-morphisms}
若存在连接 $\operatorname{tr}_1(C)$ 中两个态射的同伦，则称它们\emph{同伦}。
\VerdierSemanticAnchor{ch1-3-1-4-equivalence-relation}
“$(g,h,a)$ 与 $(g',h',a')$ 同伦”是一个与态射的直和及复合相容的等价关系。
\VerdierSemanticAnchor{ch1-3-1-4-tr1-homotopy-category}
因此可引入范畴 $\operatorname{Tr}_1(C)$：其对象是
$\operatorname{tr}_1(C)$ 的对象，其态射是 $\operatorname{tr}_1(C)$ 的态射模同伦。

\VerdierSemanticAnchor{ch1-3-1-5}
\noindent\textbf{命题 3.1.5.}\quad
\emph{设
$D,D'\colon (X\xrightarrow{f}Y)\rightrightarrows
(X'\xrightarrow{f'}Y')$ 是 $\operatorname{tr}_1(C)$ 的两个态射。若
$D$ 与 $D'$ 在 $\operatorname{tr}_1(C)$ 中同伦，则 $c(D)$ 与 $c(D')$
在 $\operatorname{comp}(C)$ 中同伦。}

\VerdierSemanticAnchor{ch1-3-1-5-proof-setup}
令 $D=(g,h,a)$、$D'=(g',h',a')$。假定 $D$ 与 $D'$ 在
$\operatorname{tr}_1(C)$ 中同伦。于是存在同伦 $b$ 与 $b'$，使得：

\VerdierSemanticAnchor{ch1-3-1-5-eq1}
\begin{equation}
  \begin{aligned}
    g'-g&=bd+db,\\
    h'-h&=b'd+db'.
  \end{aligned}
  \tag{3.1.5.1}
\end{equation}

（为简化记号，此计算省略了表示各分量次数的上标。）
\VerdierSemanticAnchor{ch1-3-1-5-second-order-homotopy}
此外，存在二阶同伦 $\alpha$，使得
\begin{equation*}
  \alpha d-d\alpha=b'f+a-(f'b+a').
\end{equation*}

\VerdierSemanticAnchor{ch1-3-1-5-cone-homotopy}
直接计算表明
\begin{equation*}
  \begin{pmatrix}
    b' & \alpha\\
    0  & -b
  \end{pmatrix}
\end{equation*}
是连接 $c(D)$ 与 $c(D')$ 的同伦。

\VerdierSemanticAnchor{ch1-3-1-6}
\noindent\textbf{推论 3.1.6.}\quad
\emph{函子}
\begin{equation*}
  c\colon\operatorname{tr}_1(C)\longrightarrow\operatorname{comp}(C)
  \qquad\text{和}\qquad
  s\colon\operatorname{comp}(C)\longrightarrow\operatorname{tr}_1(C)
\end{equation*}
\emph{与 $\operatorname{tr}_1(C)$ 和 $\operatorname{comp}(C)$ 中的同伦相容；
取商后，它们定义仍记作}

\VerdierSemanticAnchor{ch1-3-1-6-descended-functors}
\begin{equation*}
  c\colon\operatorname{Tr}_1(C)\longrightarrow K(C)
  \qquad\text{和}\qquad
  s\colon K(C)\longrightarrow\operatorname{Tr}_1(C)
\end{equation*}
\emph{的函子。}
\VerdierSemanticAnchor{ch1-3-1-6-adjunction}
\emph{函子 $c\colon\operatorname{Tr}_1(C)\longrightarrow K(C)$ 是函子
$s\colon K(C)\longrightarrow\operatorname{Tr}_1(C)$ 的左伴随。}

\VerdierSemanticAnchor{ch1-3-1-6-proof}
这由 (3.1.5) 和 (3.1.3) 形式地推出。

\VerdierSemanticAnchor{ch1-3-1-7}
\noindent\textbf{3.1.7.}\quad
\VerdierSemanticAnchor{ch1-3-1-7-s-prime}
以 $s'\colon\operatorname{comp}(C)\longrightarrow\operatorname{tr}_1(C)$
表示函子 $X\longmapsto(X\longrightarrow 0)$。
\VerdierSemanticAnchor{ch1-3-1-7-comp-translation}
函子 $cs'\colon\operatorname{comp}(C)\longrightarrow\operatorname{comp}(C)$
等于 $\operatorname{comp}(C)$ 的平移函子。
\VerdierSemanticAnchor{ch1-3-1-7-descended-s-prime}
函子 $s'$ 将 $\operatorname{comp}(C)$ 的同伦态射映为
$\operatorname{tr}_1(C)$ 的同伦态射，因而定义函子
$s'\colon K(C)\longrightarrow\operatorname{Tr}_1(C)$，仍以 $s'$ 表示。
\VerdierSemanticAnchor{ch1-3-1-7-k-translation}
函子 $cs'\colon K(C)\longrightarrow K(C)$ 等于 $K(C)$ 的平移函子。
